\documentclass[11pt]{article}
\usepackage[a4paper,margin=24mm]{geometry}
\usepackage[T1]{fontenc}
\usepackage{amsmath,amssymb,amsthm,mathtools}
\usepackage[hidelinks]{hyperref}
\allowdisplaybreaks
\setlength{\emergencystretch}{3em}
\newtheorem{theorem}{Theorem}
\newtheorem{lemma}[theorem]{Lemma}
\newtheorem{proposition}[theorem]{Proposition}
\newcommand{\C}{\mathbb C}
\newcommand{\R}{\mathbb R}
\newcommand{\tr}{\operatorname{tr}}
\newcommand{\diag}{\operatorname{diag}}
\newcommand{\norm}[1]{\left\lVert#1\right\rVert}
\title{Explicit native moment precision, inverse margins, and the original signed receiver}
\author{}
\date{20 September 2026}
\begin{document}\maketitle

This calculation supplies quantitative conditioning and precision
constants for the actual finite-moment receiver in \cite{FM} and the
deflation-aware rank-one refinement RR1--16 in \cite{RR}. Both complete
proofs were read for this derivation. Positivity is proved from the
original native envelope, rather than assumed of a numerical midpoint.
All monomial coordinates, masses, parity blocks, complex transport,
observation and conductor fibres, and signed endpoint errors remain
unchanged. The elementary quadrature context is the human work of
Zimmerling, Druskin and Simoncini \cite{ZDS}; no finite-matrix
nondeflation hypothesis is imported from it.

\section{The original measure and explicit interval floor}

Retain the stipulated quartet domain and every original constant:
\begin{equation}\label{eq:objects}
\begin{gathered}
 \rho_*=\tfrac12+\delta+i\gamma,\quad 0<\delta<\tfrac12,
 \quad\gamma>2,\quad m\ge1,\quad g=2\xi,\\
 h(s)=\prod_{z\in\{\rho_*,\bar\rho_*,1-\rho_*,1-\bar\rho_*\}}(s-z)^m,
 \quad w_h(y)=\frac{|(g/h)(1/2+iy)|^2}{2\pi},\\
 m_k=w_h^{*k},\quad k=4\ell+1\ge9,\quad
 c=k/2,\quad S=c+iy,\quad \mu_h=\int_\R w_h(y)dy,\\
 e_k=1+k(m-1),\quad q=e_k(k+1)^2,\quad
 \chi(S)=\prod_{a,b=0}^k
 [S-c-(2a-k)\delta-i(2b-k)\gamma]^{e_k}.
\end{gathered}
\end{equation}
This is the existing quartet domain, not an existence assertion for
an off-line zero. The complete original provider \cite{CAI},
CAI16--17, and its multiplicity continuation \cite{HG}, HG21, give
\begin{equation}\label{eq:envelope}
\begin{gathered}
 a_k(1+y^2)^{-M}\rho_\sigma(y)\le m_k(y)\le A_k\rho_\sigma(y),
 \quad \rho_\sigma(y)=\frac{|\Gamma(1/4+iy/2)|^2}{2\pi},\\
 c_\Gamma e^{-\alpha|y|}(1+|y|)^{-1/2}
 \le\rho_\sigma(y)\le
 C_\Gamma e^{-\alpha|y|}(1+|y|)^{-1/2},\\
 \alpha=\pi/2,\quad M=21+4m,\quad p=8m-9/2,\quad B_h=42+8m,\\
 c_\Gamma=\sqrt2e^{-7/6},\qquad C_\Gamma=\sqrt{2\sqrt5}e^{2/3},\\
 \vartheta_h=\int_{-1}^1w_h(y)dy,\qquad
 M_\alpha=\int_\R e^{\alpha y}w_h(y)dy,\\
 a_k=\frac{c_h\vartheta_h^{k-3}e^{-\alpha(k-3)}(k-2)^{-B_h}}
 {C_\Gamma2^{B_h/2}},\qquad
 A_k=\frac{C_h^{\rm tilt}}{c_\Gamma}k^{p+1}M_\alpha^{k-1}.
\end{gathered}
\end{equation}
The constants $c_h,C_h^{\rm tilt}$ are exactly the provider's TW7
lower-envelope constant and BT12 compact/tail maximum. They and the
displayed original scalar integrals are not assigned numerical values
in this note. The Gamma density is an envelope only, never a
replacement measure.

The original parity and tilted measures are
\begin{equation}\label{eq:measures}
 d\nu_0(x)=m_k(\sqrt x)x^{-1/2}dx,\quad d\nu_1=x\,d\nu_0,
 \quad s(x)=\frac{\tanh(\alpha\sqrt x)}{\sqrt x},\quad
 d\widetilde\nu_\epsilon=s\,d\nu_\epsilon\quad(\epsilon=0,1).
\end{equation}
The original positive Stieltjes identity \cite{PP,FM} is
\begin{equation}\label{eq:stieltjes}
 s(x)=\frac4\pi\sum_{r=0}^\infty\frac1{x+(2r+1)^2},\qquad
 \frac1{s(x)}=\frac2\pi+\frac4\pi\sum_{j=1}^\infty\frac{x}{x+4j^2},
 \qquad 0<s(x)\le\alpha.
\end{equation}
Evenness and $x=y^2$ give, without any mass change,
\[
 \int|E(y^2)+yO(y^2)|^2m_k(y)dy
 =\int|E|^2d\nu_0+\int|O|^2d\nu_1.
\]
In particular $\nu_0((0,\infty))=\mu_h^k$ and
$\nu_1((0,\infty))=k\eta_2\mu_h^{k-1}$, where
$\eta_2=\int y^2w_h(y)dy$. The tilted masses remain their actual
integrals, and $\widetilde\nu_1=x\widetilde\nu_0$.

\begin{lemma}[An actual native floor]\label{lem:floor}
Put
\begin{equation}\label{eq:densityfloor}
 d_\nu=\frac{a_kc_\Gamma3^{-M}e^{-\alpha\sqrt2}}
 {\sqrt2\sqrt{1+\sqrt2}},\qquad
 d_{\widetilde\nu}=\frac4{3\pi}d_\nu.
\end{equation}
For either parity, the density of $\nu_\epsilon$ on $[1,2]$ is at
least $d_\nu>0$, and that of $\widetilde\nu_\epsilon$ is at least
$d_{\widetilde\nu}>0$.
\end{lemma}
\begin{proof}
On this interval $1+x\le3$, $\sqrt x\le\sqrt2$, and
$1+\sqrt x\le1+\sqrt2$. Substitution in \eqref{eq:envelope}
and the factor $x^{-1/2}$ gives the first floor for $\nu_0$;
the additional factor $x\ge1$ gives it for $\nu_1$.
The first positive term of \eqref{eq:stieltjes} gives
$s(x)\ge4/[\pi(x+1)]\ge4/(3\pi)$, which proves the second.
No interval-restricted measure replaces an original measure: we
only lower-bound its nonnegative quadratic integral by one part.
\end{proof}

For later upper margins define, for $\epsilon=0,1$ and $r\ge0$,
\begin{equation}\label{eq:momentcaps}
 U_{\nu_\epsilon,r}=
 \frac{2A_kC_\Gamma[2(r+\epsilon)]!}
 {\alpha^{2(r+\epsilon)+1}},\qquad
 U_{\widetilde\nu_\epsilon,r}=\alpha U_{\nu_\epsilon,r}.
\end{equation}
Then $\int x^r d\lambda\le U_{\lambda,r}$ for each of these
four actual measures. Indeed the substitution in
\eqref{eq:measures} gives $\int y^{2(r+\epsilon)}m_k(y)dy$;
discard $(1+|y|)^{-1/2}\le1$ in the upper envelope and integrate
$2\int_0^\infty y^Le^{-\alpha y}dy=2L!/\alpha^{L+1}$.
The tilted bound uses $s\le\alpha$.

\section{An exact rational Gram bound in the unchanged monomials}

For $t\ge0$, set $d=t+1$ and
\begin{equation}\label{eq:J}
 J_t=\left[\frac{2^{i+j+1}-1}{i+j+1}\right]_{i,j=0}^t,
 \quad D_t^{\rm int}=\det J_t,
 \quad T_t^{\rm int}=\tr J_t,
 \quad g_t=\frac{D_t^{\rm int}}{(T_t^{\rm int})^t}.
\end{equation}
All these are explicitly evaluated rational expressions in $t$.
The convention is $g_0=1$.

\begin{lemma}\label{lem:J}
The matrix $J_t$ is positive definite and
\begin{equation}\label{eq:Jfloor}
 J_t\succeq g_tI,\qquad
 D_t^{\rm int}=\prod_{j=0}^t
 \frac{(j!)^4}{(2j)!(2j+1)!}
 =\prod_{j=0}^t\frac1{(2j+1)\binom{2j}j^2}.
\end{equation}
In particular $g_0=1$, $g_1=1/40$, and $g_2=5/981552$.
An entirely elementary lower estimate valid for all $t$ is
\begin{equation}\label{eq:gsize}
 g_t\ge(3/8)^t(2t+1)^{-(t+1)}4^{-(2t^2+t)}.
\end{equation}
\end{lemma}
\begin{proof}
For a nonzero coefficient column $z$,
$z^*J_tz=\int_1^2|\sum_{j=0}^tz_jx^j|^2dx>0$;
a nonzero polynomial cannot vanish on an interval. If the eigenvalues
are $0<\lambda_1\le\cdots\le\lambda_d$, then
$\det J_t=\lambda_1\prod_{j=2}^d\lambda_j
\le\lambda_1(\tr J_t)^{d-1}$, proving the floor.

For the determinant identity alone put $x=1+u$. The coefficient
matrix $B_{lj}=\binom jl$ for $l\le j$ has determinant one, and
$J_t=B^*H_t^{\rm Hil}B$, where
$H_t^{\rm Hil}=[1/(i+j+1)]$. This auxiliary determinant identity
does not replace any native coefficient coordinates. The Cauchy
determinant identity yields
\[
 \det H_t^{\rm Hil}
 =\frac{\prod_{0\le i<j\le t}(j-i)^2}
 {\prod_{i,j=0}^t(i+j+1)}.
\]
For completeness, multiply the determinant of $[1/(u_i+v_j)]$
by $\prod_{i,j}(u_i+v_j)$. The resulting polynomial is alternating
in each list, hence divisible by the two Vandermonde products;
its total degree equals their combined degree. The quotient is a
constant, equal to one by the one-dimensional case and induction
on the residue at $u_t=-v_t$. Taking $u_i=i$, $v_j=j+1$ proves
the displayed formula. Grouping successive factors as $t$ grows
gives the product in \eqref{eq:Jfloor}. Direct substitution gives
the three small values. Finally $\binom{2j}j\le4^j$ implies
$D_t^{\rm int}\ge(2t+1)^{-(t+1)}4^{-t(t+1)}$, while
$T_t^{\rm int}\le2\sum_{i=0}^t4^i\le(8/3)4^t$.
Division proves \eqref{eq:gsize}.
\end{proof}

\section{Every required weighted native inverse has an explicit margin}

Fix one of the actual measures $\lambda$ in \eqref{eq:measures},
write $d_\lambda=d_\nu$ or $d_{\widetilde\nu}$, and put
$\mu_r=\int x^r d\lambda$. For a real polynomial
$w(x)=\sum_{l=0}^Lw_lx^l$ with $w_l\ge0$ and $w(1)>0$, set
\begin{equation}\label{eq:W}
 W_t(w)=\left[\sum_{l=0}^Lw_l\mu_{i+j+l}\right]_{i,j=0}^t,
 \quad h_{\lambda,t}=d_\lambda g_t,
 \quad \ell_t(w)=w(1)h_{\lambda,t},\quad
 U_t(w)=\sum_{i=0}^t\sum_{l=0}^Lw_lU_{\lambda,2i+l}.
\end{equation}
On $[1,2]$, $w(x)\ge w(1)$, so the proved native density floor
and \eqref{eq:Jfloor} give the full quantitative bounds
\begin{equation}\label{eq:Wfloor}
 \boxed{\ell_t(w)I\preceq W_t(w)\preceq U_t(w)I.}
\end{equation}
The upper bound follows from positivity and the trace cap. In
particular the moment systems actually used in \cite{FM} have
the following constants, not merely existential eigenvalues:
\[
\begin{array}{c|c|c}
 w(x)&\text{matrix}&\ell_t(w)/h_{\lambda,t}\ \ (=\sum_lw_l)\\ \hline
 x^l&H_t^{(l)}=[\mu_{i+j+l}]&1\\
 x+a&H_t^{(1)}+aH_t&1+a\\
 x(x+a)&H_t^{(2)}+aH_t^{(1)}&1+a\\
 x(x+a)(x+b)&H_t^{(3)}+(a+b)H_t^{(2)}+abH_t^{(1)}&(1+a)(1+b)
\end{array}
\]
Here $a,b>0$; in the original calculation $a=4j^2$ and
$b=(2r+1)^2$. These formulas apply separately to both parity
blocks, both with and without the actual factor $s$.

\begin{theorem}[A pole-independent moment budget]\label{thm:budget}
Suppose each needed scalar moment has a real certified midpoint
$\widehat\mu_r$ and absolute error at most $\varepsilon$.
Let $\widehat W$ be its moment matrix. For $d=t+1$,
\begin{equation}\label{eq:deltaW}
 \norm{W-\widehat W}\le\delta_W:=d\,w(1)\varepsilon.
\end{equation}
For every $0<\theta\le1/4$, the concrete common budget
\begin{equation}\label{eq:uniformbudget}
 \boxed{\varepsilon\le\frac{\theta d_\lambda g_t}{t+1}}
\end{equation}
ensures $\delta_W\le\theta\ell_t(w)$ for every weight in
the table, simultaneously for all positive poles. In particular
\begin{equation}\label{eq:inversebracket}
 \widehat W-\delta_W I\succeq(\ell_t(w)-2\delta_W)I>0,
 \quad
 (\widehat W+\delta_W I)^{-1}\preceq W^{-1}
 \preceq(\widehat W-\delta_W I)^{-1},
\end{equation}
and
\begin{equation}\label{eq:inverseerror}
 \norm{W^{-1}-\widehat W^{-1}}
 \le\frac{\delta_W}{\ell_t(w)(\ell_t(w)-\delta_W)}.
\end{equation}
\end{theorem}
\begin{proof}
Each entry error is at most $w(1)\varepsilon$. Its maximum
absolute row sum and hence the norm of the Hermitian error are
at most $d w(1)\varepsilon$. Weyl's elementary quadratic-form
inequality gives $\widehat W\succeq(\ell_t(w)-\delta_W)I$.
This proves the first assertion in \eqref{eq:inversebracket};
the other assertions use inversion of positive definite Loewner
inequalities and
$W^{-1}-\widehat W^{-1}=W^{-1}(\widehat W-W)\widehat W^{-1}$.
The coefficient sum and the interval minimum are the same $w(1)$,
so it cancels in \eqref{eq:uniformbudget}.
\end{proof}

For different moment radii $\varepsilon_r$, a sharper completely
explicit choice is
\[
 \delta_W=\max_{0\le i\le t}\sum_{j=0}^t
                \sum_lw_l\varepsilon_{i+j+l}.
\]
For a finite list of degrees and measures, take the minimum of the
right side of \eqref{eq:uniformbudget}. This is a specified finite
positive number, not a request to test an unproved midpoint
positivity. When a common base sequence represents both parities,
use the literal shift $\mu_{1,r}=\mu_{0,r+1}$, preserving its
correlation and mass.

\section{Outward quadratic expressions and explicit absolute budgets}

Every inverse-derived matrix in \cite{FM} has the form
$F=C^*W^{-1}C$, where $C=[\mu_{i+j+l_0}]$ is $d$ by $e=n+1$,
$l_0=0$ or $1$, and $d=t+1\ge e$. Define explicit numbers
\begin{equation}\label{eq:P}
 \begin{gathered}
 c_C=\sqrt{de},\qquad
 U_C=\left(\sum_{i=0}^{d-1}\sum_{j=0}^{e-1}
                     U_{\lambda,i+j+l_0}^{2}\right)^{1/2},\\
 c_W=d\,w(1),\quad \ell=\ell_t(w),\qquad
 P_F=\frac43\left\{\frac{c_WU_C^2}{\ell^2}
                  +\frac{2U_Cc_C+c_C^2}{\ell}\right\}.
 \end{gathered}
\end{equation}
These are evaluated algebraic expressions in the original upper
caps and lower floors. For $\varepsilon\le1$ and
$c_W\varepsilon\le\ell/4$, set
$\widehat F=\widehat C^*\widehat W^{-1}\widehat C$.
Then
\begin{equation}\label{eq:Ferror}
 \norm{F-\widehat F}\le P_F\varepsilon,
 \qquad \widehat F-P_F\varepsilon I\preceq F
 \preceq\widehat F+P_F\varepsilon I.
\end{equation}
Indeed $\norm{C}\le U_C$, $\norm{C-\widehat C}\le c_C\varepsilon$,
$\norm{\widehat W^{-1}}\le4/(3\ell)$ and
$\norm{W^{-1}-\widehat W^{-1}}\le4c_W\varepsilon/(3\ell^2)$.
Insert these in
\[
 F-\widehat F=C^*(W^{-1}-\widehat W^{-1})C
              +C^*\widehat W^{-1}C-
                    \widehat C^*\widehat W^{-1}\widehat C.
\]
The second difference has norm at most
$(2U_Cc_C\varepsilon+c_C^2\varepsilon^2)4/(3\ell)$;
$\varepsilon^2\le\varepsilon$ proves the claim. Thus a prescribed
absolute matrix error $\tau>0$ is attained with the explicit budget
\begin{equation}\label{eq:budgetF}
 \varepsilon\le\min\{1,\ell/(4c_W),\tau/P_F\}.
\end{equation}
The restriction $\varepsilon\le1$ selects an absolute accuracy;
it is not a normalization of a measure or of its moments.

There is also a native lower floor for the exact $F$ itself. The
first $e$ rows of $C$ form $H_n^{(l_0)}\succeq h_{\lambda,n}I$,
and $W^{-1}\succeq U_t(w)^{-1}I$. Consequently
\begin{equation}\label{eq:Ffloor}
 F\succeq\frac{h_{\lambda,n}^2}{U_t(w)}I.
\end{equation}
Choosing $P_F\varepsilon\le h_{\lambda,n}^2/[4U_t(w)]$
makes the lower outward endpoint in \eqref{eq:Ferror} strictly
positive with at least half this floor. This includes the
quadratic-denominator matrix with $l_0=1$ and
$w=x(x+a)(x+b)$, not only the one-pole systems.

An actual floating computation of the midpoint inverse must itself
be certified. One option is exact rational inversion, since the
moment midpoints may be rational and the original poles are positive
integers. More generally, if a computed matrix $R_0$ satisfies a
certified $\norm{I-\widehat W R_0}\le r_0<1$, then
\begin{equation}\label{eq:residualinverse}
 \norm{\widehat W^{-1}-R_0}
 \le\frac{\norm{R_0}r_0}{1-r_0}.
\end{equation}
To prove this put $E=I-\widehat W R_0$ and use
$\widehat W^{-1}=R_0(I-E)^{-1}$ and its convergent geometric
series. Replacing $R_0$ by its Hermitian part does not increase
the error from the exact Hermitian inverse. Its extra quadratic
error is at most $\norm{\widehat C}^2$ times the right side of
\eqref{eq:residualinverse}, plus certified multiplication roundoff.
All subsequent scalar coefficients, square roots and logarithms are
likewise rounded outwards if evaluated numerically.

\section{A quantitative Radau margin and the exact complex phase}

For the actual tilted measure $\rho=\widetilde\nu_\epsilon$,
use the original RR notation, with $r\ge n+1$:
\begin{equation}\label{eq:RR}
 \begin{gathered}
 H=[\mu_{i+j}]_{0\le i,j<r},\quad
 D=[\mu_{i+j+1}]_{0\le i,j<r},\quad
 C=[\mu_{i+j}]_{\substack{0\le i<r\\0\le j\le n}},\quad
 e=(0,\ldots,0,1)^T,\\
 \kappa=(e^*D^{-1}e)^{-1},\quad \widetilde D=D-\kappa ee^*,
 \quad W_G=D+aH,\quad W_R=\widetilde D+aH.
 \end{gathered}
\end{equation}
Here $\widetilde D$ is semidefinite, not positive definite, and
is never inverted by itself. The attained constrained minimum
\begin{equation}\label{eq:kmin}
 \kappa=\min_{e^*z=1}z^*Dz
\end{equation}
has minimizer $\kappa D^{-1}e$; subtracting it from an arbitrary
point in the fibre proves $\widetilde D\succeq0$ and its exact
one-dimensional kernel. Set
\begin{equation}\label{eq:RRconstants}
 \begin{gathered}
 h=d_{\widetilde\nu}g_{r-1},\quad
 U_H=\sum_{i=0}^{r-1}U_{\rho,2i},\quad
 U_D=\sum_{i=0}^{r-1}U_{\rho,2i+1},\quad U_\kappa=U_{\rho,2r-1}.
 \end{gathered}
\end{equation}
We have
\begin{equation}\label{eq:radaufloor}
 H,D\succeq hI,\quad 0<\kappa\le D_{r-1,r-1}\le U_\kappa,
 \quad W_G\succeq(1+a)hI,\quad W_R\succeq ahI.
\end{equation}
Thus the apparent Schur denominator has the explicit native floor
\begin{equation}\label{eq:denominator}
 \boxed{1-\kappa e^*W_G^{-1}e\ge
                    \frac{ah}{U_D+aU_H}>0.}
\end{equation}
To check it, conjugate $W_R\succeq ahI$ by $W_G^{-1/2}$.
The resulting matrix is $I-\kappa uu^*$, $u=W_G^{-1/2}e$,
and is at least $ah(U_D+aU_H)^{-1}I$. Its least eigenvalue
is precisely the left side of \eqref{eq:denominator}.

This lower margin has a concrete moment tolerance. For moment
error $\varepsilon$ with $r\varepsilon<h$, positivity of
$\widehat D$ follows already from the interval floor, so define
the exact midpoint scalar $\widehat\kappa=(e^*\widehat D^{-1}e)^{-1}$.
The error bound $\norm{D-\widehat D}\le r\varepsilon$ gives
\[
 (1-r\varepsilon/h)D\preceq\widehat D
 \preceq(1+r\varepsilon/h)D.
\]
Minimizing on the same fibre in \eqref{eq:kmin} proves
\begin{equation}\label{eq:kerror}
 |\widehat\kappa-\kappa|\le
                  \frac{rU_\kappa}{h}\varepsilon.
\end{equation}
Therefore, for
$\widehat W_R=\widehat D+a\widehat H-\widehat\kappa ee^*$,
\begin{equation}\label{eq:radauprecision}
 \norm{W_R-\widehat W_R}\le c_R\varepsilon,
 \quad c_R=r(1+a)+\frac{rU_\kappa}{h},\quad
 \varepsilon\le\min\{1,h/(4r),ah/(4c_R)\}
\end{equation}
guarantees the same positive outward inverse brackets as
\eqref{eq:inversebracket}, with $\ell=ah$, $c_W=c_R$.
Formula \eqref{eq:P} with those constants certifies
$F^+=C^*W_R^{-1}C$. Its Gauss counterpart
$F^-=C^*W_G^{-1}C$ uses $\ell=(1+a)h$, $c_W=r(1+a)$.
This proves a finite precision budget even for the deflated Radau
system; no generic numerical positivity assumption was used.

For clarity, the rank-one phase assertion used in this receiver
also survives the outward calculation. Let $p_r$ be the actual
monic polynomial orthogonal for $\rho$ and let $z=v_n(-a)$.
Orthogonality of the Gauss residual gives, for every
$f(x)=v_n(x)^Tu$,
\[
 f-(x+a)q=-q_{r-1}p_r,\qquad q_{r-1}=-f(-a)/p_r(-a),
 \quad[q]=W_G^{-1}Cu.
\]
The value $p_r(-a)$ cannot vanish, since division by $x+a$
would contradict positivity of $\int(x+a)|p_r/(x+a)|^2d\rho$.
Expanding the variational square yields
\begin{equation}\label{eq:phase}
 F_n(a)=\int\frac{v_nv_n^*}{x+a}d\rho
       =F^-+\lambda zz^*,\qquad
 \lambda=\frac1{p_r(-a)^2}\int\frac{p_r(x)^2}{x+a}d\rho\ge0.
\end{equation}
The rank-one inverse identity and
$e^*W_G^{-1}C=-z^*/p_r(-a)$ give
\[
 F^+-F^-=\gamma zz^*,\qquad
 \gamma=\frac{\kappa}{p_r(-a)^2(1-\kappa e^*W_G^{-1}e)}.
\]
For the upper sign, let $t_{r-1}$ be monic orthogonal for $x\rho$.
Then $\int x|t_{r-1}|^2d\rho=\kappa$, and the Radau equations
give $f-(x+a)\widetilde q=-\widetilde q_{r-1}xt_{r-1}$ with
$\widetilde q_{r-1}=f(-a)/[a t_{r-1}(-a)]$.
Divide $\bar f(x)-\bar f(-a)$, then
$t_{r-1}(x)-t_{r-1}(-a)$, by $x+a$ and use orthogonality
against $x\rho$. The complete difference is
\[
 u^*(F^+-F_n)u=
 \frac{|f(-a)|^2}{a t_{r-1}(-a)^2}
                   \int\frac{x t_{r-1}(x)^2}{x+a}d\rho\ge0.
\]
The divided differences vanish for constants, so $r=1$ is included.
Thus $0\le\lambda\le\gamma$ with its original mass.

If \eqref{eq:Ferror} certifies midpoint errors $\eta_\pm$,
one may evaluate $\gamma$ stably without polynomial division:
\begin{equation}\label{eq:gammaout}
 \widehat\gamma=(\widehat F^+-\widehat F^-)_{00},\qquad
 \max\{0,\widehat\gamma-\eta_+-\eta_-\}\le\gamma
 \le\widehat\gamma+\eta_++\eta_-=: \gamma^{\rm up}.
\end{equation}
The exact matrix description is
\begin{equation}\label{eq:phaseout}
 F_n=\widehat F^-+E^-+\lambda zz^*,\qquad
 \norm{E^-}\le\eta_-,\quad E^-=(E^-)^*,\quad
 0\le\lambda\le\gamma^{\rm up}.
\end{equation}
In particular its complex cross error is
$u^*E^-v+\lambda\overline{f_u(-a)}f_v(-a)$, not an
independently assigned phase. The same $E^-$ and scalar $\lambda$
serve every entry. This distinction is retained by every congruence.

\section{Finite tilted and resolvent enclosures from these budgets}

Here are the precise outward substitutions in the previous native
formulas. They also specify the cost of every finite moment error.
For the chosen $\lambda$, target $n$, and auxiliary $t\ge n$,
put $H_n=[\mu_{i+j}]$, $C=[\mu_{i+j}]_{0\le i\le t,0\le j\le n}$,
$D_0=[\mu_{i+j+1}]_{0\le i\le t,0\le j\le n}$ and
\[
 L=C^*(H_t^{(1)}+aH_t)^{-1}C,\qquad
 V=D_0^*(H_t^{(2)}+aH_t^{(1)})^{-1}D_0.
\]
Complete squares with weights $x+a$ and $x(x+a)$ respectively
give
\begin{equation}\label{eq:variational}
 L\preceq R_n(a)\preceq(H_n-V)/a,\quad
 V\preceq B_n(a)\preceq H_n-aL,
\end{equation}
where $R_n=\int vv^*/(x+a)d\lambda$,
$B_n=\int xvv^*/(x+a)d\lambda$. Explicitly the two lower
errors on a vector $u$ are the minima over degree-$t$ polynomials
$q$ of $\int(x+a)|f/(x+a)-q|^2d\lambda$ and
$\int x(x+a)|f/(x+a)-q|^2d\lambda$, respectively.
Subtracting each lower bound from $H_n=aR_n+B_n$ proves the
upper bounds. Thus no sign in \eqref{eq:variational} is assumed.

Let $\eta_H=(n+1)\varepsilon$ and $\eta_L=P_L\varepsilon$,
$\eta_V=P_V\varepsilon$. Certified finite endpoints are
\begin{equation}\label{eq:outvar}
 \begin{aligned}
 R^{\rm lo}&=\widehat L-\eta_LI,&
 R^{\rm up}&=(\widehat H_n-\widehat V+(\eta_H+\eta_V)I)/a,\\
 B^{\rm lo}&=\widehat V-\eta_VI,&
 B^{\rm up}&=\widehat H_n-a\widehat L+(\eta_H+a\eta_L)I.
 \end{aligned}
\end{equation}
For Radau replace these by
$B^{\rm lo}=\widehat H_n-a\widehat F^+-(\eta_H+a\eta_+)I$
and
$B^{\rm up}=\widehat H_n-a\widehat F^-+(\eta_H+a\eta_-)I$.
They enclose the exact $H_n-aF^+$ and $H_n-aF^-$ endpoints.

When only untilted moments are used, keep the exact mixed-pole
integral against $\nu_\epsilon$:
\[
 T_n(a,b)=\int\frac{xvv^*}{(x+a)(x+b)}d\nu_\epsilon,
 \quad T^- =D_0^*W_t(x(x+a)(x+b))^{-1}D_0.
\]
Completing the square in the weight $x(x+a)(x+b)$ for the
function $f/[(x+a)(x+b)]$ proves $T^-\preceq T_n$.
Its lower outward matrix is $\widehat T^--P_T\varepsilon I$.
The exact partial fraction identity gives the upper matrix
\begin{equation}\label{eq:partial}
 T^{\rm up}=
 \begin{cases}
 [bR^{\rm up}(b)-aR^{\rm lo}(a)]/(b-a),&b>a,\\
 [aR^{\rm up}(a)-bR^{\rm lo}(b)]/(a-b),&a>b.
 \end{cases}
\end{equation}
Indeed $T_n=[bR_n(b)-aR_n(a)]/(b-a)$; the indicated upper
and lower signs follow from the sign of each coefficient. The
original even-square and odd-square poles are never equal.

For explicit scalar precision accounting, the centers for the
exact finite $R^-$ and $R^+$ expressions have error coefficients
$\eta_{R^-}(a)=\eta_L(a)$ and
$\eta_{R^+}(a)=(\eta_H+\eta_V(a))/a$. Consequently the
center for the finite upper expression in \eqref{eq:partial}
has radius
\[
 \eta_{T^+}(a,b)=
 \begin{cases}
 [b\eta_{R^+}(b)+a\eta_{R^-}(a)]/(b-a),&b>a,\\
 [a\eta_{R^+}(a)+b\eta_{R^-}(b)]/(a-b),&a>b,
 \end{cases}
 \qquad \eta_{T^-}=P_T\varepsilon.
\]
These are the sums of absolute signed coefficients, not
unsupported subtraction of error bounds.

For $b_r=(2r+1)^2$ and any integer $R\ge1$ put
\[
 \theta_R=\alpha-\frac4\pi\sum_{r=0}^{R-1}\frac1{b_r}>0,
 \qquad\theta_R\le\frac2{\pi(2R-1)}.
\]
The first equality uses the value of the positive series at zero;
the bound compares the decreasing omitted terms with the integral
of $(2u+1)^{-2}$ from $R-1$ to infinity. If
$H^{\rm up}=\widehat H_n+\eta_HI$, then the actual tilted
matrices are enclosed by
\begin{equation}\label{eq:tiltedout}
 \begin{aligned}
 \frac4\pi\sum_{r<R}R^{\rm lo}(b_r)
 &\preceq \int vv^*s\,d\nu_\epsilon
 \preceq\frac4\pi\sum_{r<R}R^{\rm up}(b_r)+\theta_RH^{\rm up},\\
 \frac4\pi\sum_{r<R}T^{\rm lo}(a,b_r)
 &\preceq \int\frac{xvv^*s}{x+a}d\nu_\epsilon
 \preceq\frac4\pi\sum_{r<R}T^{\rm up}(a,b_r)+\theta_RH^{\rm up}.
 \end{aligned}
\end{equation}
The omitted positive scalar kernels are at most $\theta_R$,
since $x/(x+a)\le1$. Integration proves every remainder in
\eqref{eq:tiltedout}. These are finite original-moment enclosures
with all scalar tails and signed coefficients retained.

The four finite centers on the two lines of \eqref{eq:tiltedout}
therefore have explicit radii
\[
 \begin{array}{ll}
 \eta_{K^-}=\dfrac4\pi\sum_{r<R}\eta_{R^-}(b_r),&
 \eta_{K^+}=\dfrac4\pi\sum_{r<R}\eta_{R^+}(b_r)+\theta_R\eta_H,\\[3pt]
 \eta_{\widetilde B^-}=\dfrac4\pi\sum_{r<R}\eta_{T^-}(a,b_r),&
 \eta_{\widetilde B^+}=\dfrac4\pi\sum_{r<R}\eta_{T^+}(a,b_r)
                                      +\theta_R\eta_H.
 \end{array}
\]
Every coefficient here is the displayed finite expression in
original constants, poles and degrees; no unknown inverse norm
is left inside a precision request.

\section{Both parity blocks, outward source bounds, and direct precision}

At an original degree $N\ge q-1$ put
$n_0=\lfloor N/2\rfloor$, $n_1=\lfloor(N-1)/2\rfloor$, and
\[
 H_N=\diag\left([\int x^{i+j}d\nu_0]_{0\le i,j\le n_0},
                [\int x^{i+j}d\nu_1]_{0\le i,j\le n_1}\right).
\]
Let $K$ and $K_x$ be the corresponding actual tilted moment
blocks with no shift and one shift. The original reciprocal series
in \eqref{eq:stieltjes}, the one-pole bounds above, and
$\sum_{j>J}j^{-2}\le1/J$ give
\begin{equation}\label{eq:Gexact}
 G^-:=\frac2\pi K+\frac4\pi\sum_{j=1}^JB_j^-
 \preceq H_N\preceq
 G^+:=\frac2\pi K+\frac4\pi\sum_{j=1}^JB_j^+
                        +\frac1{\pi J}K_x.
\end{equation}
Here $B_j^\pm$ may be the exact finite variational or Radau
endpoints in both blocks. They are positive semidefinite: for
Radau, $W_R\succeq aH$ gives
$aC^*W_R^{-1}C\preceq C^*H^{-1}C=H_n$; the variational
lower endpoint is manifestly positive. Therefore
\begin{equation}\label{eq:gsource}
 G^-\succeq g_N^*I,\qquad
 g_N^*=\frac2\pi d_{\widetilde\nu}
                         \min\{g_{n_0},g_{n_1}\}>0.
\end{equation}
Apply the moment budgets to every displayed finite matrix. If
its exact endpoint has a center $\widehat G^\pm$ with certified
norm error $\eta_\pm$, set
\begin{equation}\label{eq:Gout}
 G^{\rm lo}=\widehat G^--\eta_-I,\qquad
 G^{\rm up}=\widehat G^++\eta_+I.
\end{equation}
For the direct tilted route one may take, explicitly,
\[
 \eta_-\le\frac2\pi\eta_K+\frac4\pi\sum_{j=1}^J\eta_{B_j^-},
 \quad
 \eta_+\le\frac2\pi\eta_K+\frac4\pi\sum_{j=1}^J\eta_{B_j^+}
                         +\frac1{\pi J}\eta_{K_x}.
\]
Each $\eta$ on the right is the displayed coefficient times
$\varepsilon$, or the maximum of the two parity coefficients.
Thus these give finite known coefficients $P_\pm$ with
$\eta_\pm\le P_\pm\varepsilon$; choose
$\varepsilon\le g_N^*/(4P_-)$ as well as all local budgets.
Then $G^{\rm lo}\succeq(g_N^*-2\eta_-)I\succeq g_N^*I/2$,
and
\begin{equation}\label{eq:sourcecert}
 0<G^{\rm lo}\preceq H_N\preceq G^{\rm up},\quad
 \Delta G=G^{\rm up}-G^{\rm lo}\succeq0,\quad
 G^{\rm up}\preceq(1+\eta_N)G^{\rm lo},\quad
 \eta_N=\frac{\norm{\Delta G}}{g_N^*-2\eta_-}.
\end{equation}
The exact norm may be replaced by a certified Frobenius or row-sum
upper bound. The original-moment-only route \eqref{eq:tiltedout}
uses the same signed scalar error addition, with the actual native
outer tail $X_H/(2J)$ in place of $K_x/(\pi J)$, where $X_H$
is the once-shifted untilted direct sum. This follows from
$s\le\pi/2$, with no substituted measure. In that route the
exact finite lower approximation to $K$ is smaller than $K$,
so its positivity floor must also be replaced. Keep its first
pole $b_0=1$ and let $t_{\epsilon,0}$ be that pole's actual
auxiliary degree. Formula \eqref{eq:Ffloor} gives the explicit
valid replacement
\[
 g_{N,\mathrm{untilted}}^*
 =\frac8{\pi^2}\min_{\epsilon=0,1}
 \frac{(d_\nu g_{n_\epsilon})^2}
 {U_{t_{\epsilon,0}}^{(\nu_\epsilon)}(x+1)}>0.
\]
All other finite lower summands are positive. Use this floor,
the literal signed coefficients from \eqref{eq:partial}, and
the tails in \eqref{eq:tiltedout} in
\eqref{eq:Gout}--\eqref{eq:sourcecert}; the same proved
$\eta_-\le g_{N,\mathrm{untilted}}^*/4$ budget certifies the
lower source endpoint. No positivity floor for the exact tilted
Gram is mistakenly assigned to its finite lower approximation.

There is a stronger elementary route when the original moments
through degree $2n_\epsilon$ are supplied: they already give
$H_N$ directly. There is no need to reconstruct it by resolvents.
Define
\begin{equation}\label{eq:directconstants}
 h_N=d_\nu\min\{g_{n_0},g_{n_1}\},\qquad
 d_N^{\rm src}=\max\{n_0+1,n_1+1\},\qquad
 \delta_N=d_N^{\rm src}\varepsilon.
\end{equation}
Then $H_N\succeq h_NI$ and $\norm{H_N-\widehat H_N}\le\delta_N$.
For $\delta_N\le\theta_Nh_N$, $0<\theta_N\le1/4$, put
\begin{equation}\label{eq:directout}
 H_N^{\rm lo}=\widehat H_N-\delta_NI,\quad
 H_N^{\rm up}=\widehat H_N+\delta_NI.
\end{equation}
They are positive outward endpoints and obey
\begin{equation}\label{eq:directrelative}
 H_N^{\rm up}\preceq\frac1{1-2\theta_N}H_N^{\rm lo}.
\end{equation}
Indeed $H_N^{\rm lo}\succeq(h_N-2\delta_N)I$ and the two
endpoints differ by $2\delta_NI$, giving factor
$1+2\delta_N/(h_N-2\delta_N)\le(1-2\theta_N)^{-1}$.
This direct original-moment bound, unlike a still-to-be-selected
resolvent depth, supplies an unconditional finite entrywise
precision budget for any prescribed original four-endpoint error.

\section{Exact complex transport, inverse evaluation, and full fibres}

Retain the full original coefficient map, including all phases:
\begin{equation}\label{eq:transport}
 \begin{gathered}
 (T_N)_{rj}=\binom jr c^{j-r}i^r,\quad
 (T_N^{-1})_{rj}=\binom jr(-c)^{j-r}i^{-j}\quad(r\le j),\\
 L_N=\Pi_NT_N,\quad H_S=L_N^*H_NL_N,\quad
 \det T_N=i^{N(N+1)/2},\quad
 \det L_N=(\det\Pi_N)i^{N(N+1)/2},\\
 A=\Lambda J_{S,N}L_N^{-1}.
 \end{gathered}
\end{equation}
Other triangular entries vanish and $\Pi_N$ orders even powers
before odd powers. Expanding $(c+iy)^j$ proves the forward map;
expanding $y^j=i^{-j}(S-c)^j$ proves its inverse. The original
physical Taylor unit remains in $J_{S,N}$, and $A$ is onto the
same original observation image, not replaced by its norm.

For any positive source matrix $D$ in these parity coordinates,
\begin{equation}\label{eq:Q}
 Q_A(D)=(AD^{-1}A^*)^{-1},\quad
 \min_{Az=u}z^*Dz=u^*Q_A(D)u,\quad
 z_{\min}=D^{-1}A^*Q_A(D)u.
\end{equation}
The displayed minimizer is in the fibre and is $D$-orthogonal to
$\ker A$; subtraction and expansion of the square prove both
identities. Its original $S$ lift is $L_N^{-1}z_{\min}$.
For either outward source pair above, the same-fibre minimum gives
\begin{equation}\label{eq:Qbounds}
 Q^-:=Q_A(D^{\rm lo})\preceq Q_A(H_N)\preceq
 Q^+:=Q_A(D^{\rm up}).
\end{equation}
The source relative factor is preserved exactly. Any original
specified source subspace uses its same inclusion and restricted
fibre, never a silently enlarged polynomial source.

No further unverified inverse is hidden in \eqref{eq:Qbounds}.
Here is an explicit evaluation margin. For the actual onto $A$
with $d>0$ rows put
\begin{equation}\label{eq:Afloor}
 a_A=\frac{\det(AA^*)}{[\tr(AA^*)]^{d-1}}>0,
 \quad b_A=\tr(AA^*).
\end{equation}
The determinant/trace eigenvalue proof in Lemma \ref{lem:J}
gives $a_AI\preceq AA^*\preceq b_AI$. This is the literal
original map, including its jet and complex coefficients. If
$lI\preceq D\preceq UI$, then
\begin{equation}\label{eq:obsmargins}
 \frac{a_A}{U}I\preceq AD^{-1}A^*\preceq\frac{b_A}{l}I,
\qquad \frac l{b_A}I\preceq Q_A(D)\preceq\frac U{a_A}I.
\end{equation}
These margins have already proved native choices, so the
inequality is not an additional positivity premise. For the direct
moment endpoints take, simultaneously,
\[
 l=h_N-2\delta_N,\qquad
 U=\sum_{\epsilon=0}^1\sum_{i=0}^{n_\epsilon}
                 U_{\nu_\epsilon,2i}+2\delta_N.
\]
The exact native source has norm at most the displayed trace cap,
and each endpoint is within $2\delta_N$ of it. For the finite
resolvent endpoints take $l=g_N^*-2\eta_-$ (or its displayed
untilted replacement) and $U=\tr G^{\rm up}$, increased by any
certified arithmetic error in that trace. Positivity and the order
in \eqref{eq:sourcecert} prove these common endpoint bounds.
Thus a certified numerical center for $D^{-1}$ with error $r_D$
gives an observation-center error at most $b_Ar_D$, before adding
its own certified multiplication roundoff. Requiring that total
error $\delta_B\le a_A/(4U)$ gives the outward observation
inverse by \eqref{eq:inversebracket}, with $\ell=a_A/U$.
One can meet this stated precision using exact arithmetic or
\eqref{eq:residualinverse}. It is not inferred from a midpoint
determinant. If $A$ itself is evaluated numerically with error
$\delta_A$ and midpoint $\widehat A$, its additional product
error against a source inverse of norm at most $1/l$ is at most
$(2\norm{A}\delta_A+\delta_A^2)/l$ (or replace $\norm{A}$
by $\norm{\widehat A}+\delta_A$). Its original coefficient
enclosures are then separate inputs; this note does not manufacture
them. The determinant in \eqref{eq:Afloor} is an exact formula
in the original onto map, not a universal lower bound independent
of that map. The source-relative tolerance below does not require
its numerical evaluation.

In particular, if both the source inverse and the map are
approximated, the combined product radius, before multiplication
roundoff, is explicitly
\[
 \norm{AD^{-1}A^*-\widehat A R_0\widehat A^*}
 \le \frac{2\norm{A}\delta_A+\delta_A^2}{l}
                  +\norm{\widehat A}^{\,2}r_D.
\]
This follows by first changing $A$ against the exact inverse,
then changing the inverse against $\widehat A$. It retains the
mixed effect of simultaneous map and inverse errors rather than
incorrectly using the coefficient $b_A$ for both centers.

\section{Every complex cross term and every endpoint sign}

For any fixed original observation columns $U=(u_1,u_2,w)$,
the entire $3$ by $3$ Gram is bounded by
\[
 U^*Q^-U\preceq U^*Q_A(H_N)U\preceq U^*Q^+U.
\]
Its actual common remainder $E=U^*(Q_A(H_N)-Q^-)U$ satisfies
$0\preceq E\preceq U^*(Q^+-Q^-)U$; in particular
\begin{equation}\label{eq:cross}
 |u^*(Q_A(H_N)-Q^-)v|^2
 \le(u^*\Delta Q u)(v^*\Delta Q v),\qquad\Delta Q=Q^+-Q^-.
\end{equation}
This follows by Cauchy--Schwarz in the positive remainder. It
retains the single PSD matrix, including every principal minor and
the full determinant constraint, not three independently chosen
complex errors. For example, if $E_{ii}=\alpha_i$ and
$E_{ij}=z_{ij}$ for $i<j$, its constraints include $\alpha_i\ge0$,
$|z_{ij}|^2\le\alpha_i\alpha_j$, and the full determinant
\[
 \alpha_1\alpha_2\alpha_3
 +2\Re(z_{12}z_{23}\overline{z_{13}})
 -\alpha_1|z_{23}|^2-\alpha_2|z_{13}|^2-\alpha_3|z_{12}|^2
 \ge0.
\]
Expansion of the determinant proves this expression; the same
principal-minor constraints apply to $U^*\Delta Q U-E$.
In the original currents set
$u_i=\Lambda a_i$, $a_1=b_N$, $a_2=b_{N+1}$, and
$w=\Lambda v_\lambda$ with its original terminal class and
$\omega_N>0$. Put $B_i=u_i^*Q_A(H_N)w$, $B_i^-=u_i^*Q^-w$,
$\epsilon_i=B_i-B_i^-$ and
$e_i=[(u_i^*\Delta Q u_i)(w^*\Delta Q w)]^{1/2}$.
The complete signed product is
\begin{equation}\label{eq:current}
 \frac{2}{\omega_N}\Re(\overline B_1B_2)
 =\frac{2}{\omega_N}\Re\{
 \overline{B_1^-}B_2^-+\overline{B_1^-}\epsilon_2
 +\overline{\epsilon_1}B_2^-+\overline{\epsilon_1}\epsilon_2\}.
\end{equation}
Its difference from the first term has absolute value at most
$2(|B_1^-|e_2+|B_2^-|e_1+e_1e_2)/\omega_N$.
No cross term, conjugation, remainder, or phase is assigned zero.

For the actual image dimension $d_N$, define
$e_N=\log\det Q_N^+-\log\det Q_N^-\ge0$ in the unchanged
frame. A source relative factor $r_N\ge1$ gives
\begin{equation}\label{eq:logbound}
 0\le\log\det Q_N-\log\det Q_N^-\le e_N\le d_N\log r_N.
\end{equation}
The proof conjugates by $(Q_N^-)^{-1/2}$ and multiplies the
eigenvalues in $[1,r_N]$. This auxiliary computation does not
change the specified determinant or frame. The empty determinant
is one. At exactly the original window
$(N_1,N_2,N_3,N_4)=(q-1,q,2q-1,2q)$ let
\[
 F=\log\frac{\det Q_{N_1}\det Q_{N_2}}
                   {\det Q_{N_3}\det Q_{N_4}},\quad
 F^-=\log\frac{\det Q_{N_1}^-\det Q_{N_2}^-}
                     {\det Q_{N_3}^-\det Q_{N_4}^-}.
\]
The complete signed enclosure is
\begin{equation}\label{eq:four}
 \boxed{-e_{N_3}-e_{N_4}\le F-F^-
                         \le e_{N_1}+e_{N_2}.}
\end{equation}
Equivalently the lower ratio uses lower numerator matrices and
upper denominator matrices, and the upper ratio reverses these.

For a prescribed absolute four-endpoint tolerance $t_*>0$, the
direct actual-moment route supplies the explicit endpoint budgets
\begin{equation}\label{eq:finalbudget}
 \boxed{\theta_N=\min\{1/4,t_*/(8d_N)\},\qquad
 \varepsilon_N\le
 \frac{\theta_Nd_\nu\min(g_{n_0},g_{n_1})}
      {\max(n_0+1,n_1+1)}\quad(d_N>0).}
\end{equation}
Indeed \eqref{eq:directrelative} gives
$e_N\le-d_N\log(1-2\theta_N)\le4d_N\theta_N\le t_*/2$;
the middle inequality follows by integrating
$(1-u)^{-1}\le2$ over $0\le u\le2\theta_N\le1/2$.
If the minimum selects $1/4$, then $t_*\ge2d_N$ and the last
inequality still holds. Equation \eqref{eq:four} now proves
$|F-F^-|\le t_*$. Endpoints of rank zero have zero determinant
error and require no such budget. Taking the minimum of the four
positive budgets supplies one compatible shared moment precision.
No observation conditioning hypothesis is needed for this exact
relative result, although numerical evaluation of the endpoint
matrices still follows \eqref{eq:obsmargins}.

If numerical log determinants are returned, retain their outward
intervals too. Write $l_i^-\le\log\det Q_i^-\le u_i^-$ and
$l_i^+\le\log\det Q_i^+\le u_i^+$. The fully outward real
return interval is
\[
 [\,l_1^-+l_2^--u_3^+-u_4^+,\quad
       u_1^++u_2^+-l_3^--l_4^-\,].
\]
For an absolute error target $t_*$ on a reported approximation
to $F^-$, use \eqref{eq:finalbudget} with $t_*/2$ instead of
$t_*$ and certify each of the four numerical lower-endpoint log
values to absolute error at most $t_*/8$. The signed sum then
has numerical error at most $t_*/2$, and the total return error
is at most $t_*$. This separates, and includes, the actual moment
error and the finite arithmetic error.

\section{The same conductor and its unchanged signed-state kernel}

On the already established original conductor domain \cite{TR},
retain $C=\bar C\Lambda$ with $\bar C$ onto its actual image.
The exact conductor metric and outward endpoints are
\begin{equation}\label{eq:conductor}
 T_N=(\bar C Q_N^{-1}\bar C^*)^{-1},\quad
 T_N^\pm=(\bar C(Q_N^\pm)^{-1}\bar C^*)^{-1}
       =[(\bar CA)(D^\pm)^{-1}(\bar CA)^*]^{-1}.
\end{equation}
Substitution proves the last equality; minimizing on the full same
fibre $\bar Cu=v$ proves $T_N^-\preceq T_N\preceq T_N^+$
with the identical source-relative factor. This includes the entire
conductor kernel. Its numerical inverse floor is
\eqref{eq:obsmargins} with the literal map $\bar CA$.
Keep the original cubic inclusion $I_W$, coordinate map $\Psi$,
and signed-state map $B_{\mathcal J}=O(I-D_{\mathcal J})O^{-1}E$.
For the original label vector $\alpha_0$ the outward signed-state
energies are precisely
\[
 \alpha_0^*B_{\mathcal J}^*\Psi^*I_W^*T_N^-I_W\Psi
               B_{\mathcal J}\alpha_0
 \le\alpha_0^*B_{\mathcal J}^*\Psi^*I_W^*T_NI_W\Psi
               B_{\mathcal J}\alpha_0
 \le\alpha_0^*B_{\mathcal J}^*\Psi^*I_W^*T_N^+I_W\Psi
               B_{\mathcal J}\alpha_0.
\]
These are literal congruences of \eqref{eq:conductor}. The
terminal-class residual of the original EQ20--21 is retained,
not set to zero by these positive energy bounds. Its remaining
complex pairings use \eqref{eq:cross} with the original columns.
Conductor determinants, or a specified positive definite
restriction, use \eqref{eq:four} and \eqref{eq:finalbudget} with
their actual ranks, frames and the same four endpoints.

\section{Finite native scalar inputs and what the precision costs}

The rational numbers $g_t$, all pole values, determinant phases,
dimension factors and error coefficients above are fully specified
finite expressions. No native integral has been numerically
evaluated without supplied values of the original quartet data.
The remaining scalar inputs are precisely the original
$\vartheta_h$, $M_\alpha$, TW7/BT12 constants, and finitely many
actual moments. Downward bounds for $a_k,d_\nu$ and upward bounds
for $A_k,U_{\lambda,r}$ may replace their exact values in every
budget. This preserves all inequalities; it never silently treats
an uncertified midpoint constant as a lower bound.

For the untilted route the moments can be reduced further to
original onefold scalar integrals. Write
$\eta_l=\int y^lw_h(y)dy$ and
$\zeta_L=\int y^Lm_k(y)dy$. Evenness gives $\eta_{2j+1}=0$,
and Tonelli followed by the finite multinomial formula gives
\begin{equation}\label{eq:multinomial}
 \int x^r d\nu_\epsilon=\zeta_{2(r+\epsilon)},\quad
 \zeta_L=\sum_{l_1+\cdots+l_k=L}
 \frac{L!}{l_1!\cdots l_k!}\prod_{j=1}^k\eta_{l_j}.
\end{equation}
Odd summands vanish, and all remaining factors are nonnegative.
Every monomial product in the finite expansion is absolutely
integrable by the original exponential moment. Fubini therefore
justifies the factorization even for the terms containing odd
powers, before evenness makes those terms zero.
For $L=0$ this is the actual mass $\mu_h^k$; it is not one.
The original exponential moment gives the explicit caps
\begin{equation}\label{eq:etaone}
 0\le\eta_{2j}\le(2j)!\alpha^{-2j}M_\alpha,
\end{equation}
because $\cosh(\alpha y)\ge(\alpha y)^{2j}/(2j)!$ and
$\int\cosh(\alpha y)w_h=M_\alpha$.
Suppose original moments through $\zeta_{2R}$ are wanted at
absolute error $\varepsilon$. Put
$U=\max_{0\le j\le R}(2j)!\alpha^{-2j}M_\alpha$.
If every even onefold input is enclosed with midpoint error
\begin{equation}\label{eq:onefoldbudget}
 \tau\le\min\left\{1,
 \frac{\varepsilon}{k^{2R+1}(U+1)^{k-1}}\right\},
\end{equation}
then the finite polynomial \eqref{eq:multinomial}, evaluated
exactly on those midpoints, has error at most $\varepsilon$
for every even degree through $2R$. To prove it, telescope each
$k$-fold product; its error is at most
$k\tau(U+1)^{k-1}$. The sum of all multinomial coefficients
at degree $L$ is $k^L$, and the even-only subsum is smaller.
For $L\le2R$ their product is bounded by the denominator in
\eqref{eq:onefoldbudget}. Thus the four-endpoint budget is an
explicit finite accuracy request on original onefold integrals,
not on an unidentified proxy measure.

Their infinite tails also have an original explicit certificate.
For an even degree $L$ and $Y\ge\max\{1,L/\alpha\}$,
\begin{equation}\label{eq:tail}
 \int_{|y|>Y}|y|^Lw_h(y)dy
 \le2M_\alpha Y^Le^{-\alpha Y},\quad
 \int_{|y|>Y}|y|^Lm_k(y)dy
 \le2M_\alpha^kY^Le^{-\alpha Y}.
\end{equation}
Indeed $y^Le^{-\alpha y}$ decreases for $y\ge L/\alpha$,
and evenness gives $\int e^{\alpha|y|}w_h\le2M_\alpha$.
Tonelli proves $\int e^{\alpha y}m_k=M_\alpha^k$, giving
the second bound. A tilted moment tail has the additional factor
at most $\alpha$. Repeatedly doubling $Y$ from the stated lower
threshold attains any prescribed positive tail allowance, since
$Y^Le^{-\alpha Y}\to0$. The remaining compact integrals of
the original $w_h$, or of $s(y^2)m_k(y)$ if one chooses the
tilted route, still require certified evaluation. Their numerical
values are not supplied by these tail inequalities or by the
moment-matrix calculation.

For absolute input error $2^{-p}$, \eqref{eq:uniformbudget}
requires the explicit integer
\[
 p\ge\max\left\{0,
 \left\lceil\log_2\frac{t+1}{\theta d_\lambda g_t}\right\rceil
 \right\}.
\]
There is no hidden unknown eigenvalue. Formula \eqref{eq:gsize}
shows an elementary sufficient cost with
$-\log_2g_t\le4t^2+2t+t\log_2(8/3)
 +(t+1)\log_2(2t+1)$.
These unaltered monomial-coordinate floors are deliberately
conservative and may demand very high precision at the original
large degrees. This proof establishes finite, explicit tolerances
and outward certificates, not practical low-cost evaluation of
the native integrals or a new phase/asymptotic conclusion.

The companion exact tests verify rational interval Grams,
determinants, all polynomial weights, moment perturbations, inverse
brackets, the deflated Radau margin, complex rank-one cross terms,
original triangular maps and masses, quotient/conductor fibres,
and all four determinant signs. Test measures serve only as exact
checks of the proved universal identities. They do not replace the
native measures in any statement above.

\begin{thebibliography}{9}
\bibitem{ZDS} J\"orn Zimmerling, Vladimir Druskin and Valeria Simoncini,
\emph{Monotonicity, bounds and acceleration of Block Gauss and
Gauss--Radau quadrature for computing $B^T\phi(A)B$}, original
author source of arXiv:2407.21505v3, 9 January 2025, Sections 3--5.
\url{https://arxiv.org/abs/2407.21505v3}.
The abstract-page title uses ``extrapolation''; the title here
identifies the retained version-three author source. Its complete
source was read by the integrating task. This note uses it for
human mathematical provenance, and proves its native conditioning
and perturbation claims independently.
\bibitem{CAI} Split-Zero programme, \emph{The complete original action
relative to its arithmetic total}, complete provider
\path{CAI_COMPLETE_PREREQUISITE.tex}, CAI16--17 and its original
TW7/BT12 constant definitions. The exact displayed envelope and
constants are retained above.
\bibitem{HG} Split-Zero programme, \emph{Subexponential heat cost and
the original arithmetic currents}, complete provider
\path{HEAT_GROWTH_AND_ORIGINAL_CURRENTS.tex}, HG21, arbitrary
fixed primary multiplicity and the full-measure comparison.
\bibitem{PP} Split-Zero programme, \emph{Exact native parity transport
into the Pollaczek pair and certified finite source--observation
bounds}, complete proof
\path{independent_parity_pollaczek/PARITY_POLLACZEK_NATIVE_RECEIVER.tex},
20 September 2026. The unscaled Pollaczek ratio there retains the
human source A. I. Aptekarev, G. L\'opez Lagomasino and
A. Mart\'{\i}nez--Finkelshtein, \emph{Strong asymptotics for the
Pollaczek multiple orthogonal polynomials ensembles},
arXiv:1410.1261v1 (2014). No asymptotic or replacement weight from
that paper is imported here.
\bibitem{FM} Split-Zero programme, \emph{Certified finite native moment
enclosures for the parity resolvents}, complete proof
\path{independent_resolvent_moments/NATIVE_RESOLVENT_MOMENT_ENCLOSURES.tex},
20 September 2026. Read completely as a prerequisite.
\bibitem{RR} Split-Zero programme, \emph{One scalar resolvent
uncertainty in the original polynomial observation}, complete proof
\path{NATIVE_RESOLVENT_RANK_ONE.tex}, RR1--16, 20 September 2026.
Read completely as a prerequisite; its full mass, deflation,
rank-one phase, original maps and endpoint signs are retained.
\bibitem{TR} Split-Zero programme, \emph{Native resolvent truncation at
growing degree and the original conductor minimum}, complete proof
\path{NATIVE_TRUNCATION_CONDUCTOR_RETURN.tex}, TR1--10,
20 September 2026, retaining the complete original EQ conductor
and signed-state maps, including their terminal residuals.
\end{thebibliography}
\end{document}
